\chapter{Introdu\c{c}\~ao}
\label{ch:introducao}

\section{Motiva\c{c}\~ao}

Sistemas distribu\'idos de grande escala t\^em se apoiado cada vez mais
em protocolos de consenso aproximado para equilibrar lat\^encia e
consist\^encia \cite{lamport1998part}. Este trabalho estuda como tais
protocolos se degradam sob parti\c{c}\~ao de rede.

\section{Contribui\c{c}\~oes}

As principais contribui\c{c}\~oes s\~ao:

\begin{enumerate}
  \item Um modelo formal de toler\^ancia a parti\c{c}\~oes parciais.
  \item Uma avalia\c{c}\~ao emp\'irica resumida na Tabela~\ref{tab:resultados}.
  \item Uma implementa\c{c}\~ao de refer\^encia usada para gerar a
    Figura~\ref{fig:latencia}.
\end{enumerate}

\section{Trabalhos Relacionados}

Protocolos de consenso cl\'assicos como Paxos \cite{lamport1998part} e
Raft \cite{ongaro2014search} assumem uma rede totalmente conectada.
Esta disserta\c{c}\~ao relaxa essa hip\'otese.
